\documentclass[11pt,a4paper]{article}

\usepackage[margin=1in]{geometry}
\usepackage{times}
\usepackage{amsmath,amssymb,amsthm}
\usepackage{graphicx}
\usepackage{booktabs}
\usepackage{multirow}
\usepackage{array}
\usepackage{listings}
\usepackage{xcolor}
\usepackage{algorithm}
\usepackage{algpseudocode}
\usepackage{hyperref}
\usepackage{url}
\usepackage{cite}

\hypersetup{colorlinks=true, linkcolor=blue, citecolor=blue, urlcolor=blue}

\lstset{
    basicstyle=\ttfamily\footnotesize,
    keywordstyle=\color{blue}\bfseries,
    commentstyle=\color{gray},
    stringstyle=\color{red!70!black},
    breaklines=true,
    frame=single,
    numbers=left,
    numberstyle=\tiny\color{gray},
    showstringspaces=false,
    language=Java
}

\title{\textbf{An AI-Augmented Event-Driven Reference Architecture for Pre-Authorization Decisioning in ACH Payments: Joint Fraud, Return Code, and Compliance Risk Modeling under Operational SLA Constraints}}

\author{Aravind Raghu\thanks{Corresponding author. Affiliation and contact details to be inserted.}}

\date{}

\begin{document}

\maketitle

\begin{abstract}
The Automated Clearing House (ACH) network processes more than seventy
million daily payment instructions in the United States, and the operational
risk decisions that gate each transaction --- fraud screening, return-code
forecasting, and Bank Secrecy Act / Anti-Money-Laundering (BSA/AML)
compliance review --- are traditionally serviced by independent rule-based or
machine-learned pipelines. This paper presents a Java/Spring~Boot reference
architecture for joint AI-driven pre-authorization decisioning that consolidates
these three risk surfaces behind a single shared-encoder neural network, served
through an event-driven microservice topology suitable for Tier-1 banking
deployment. We make four concrete contributions: (i) a complete reference
architecture for AI-augmented ACH decisioning under sub-100~$\mu$s
single-transaction inference SLAs; (ii) an empirical study showing that
multi-task representation sharing achieves area-under-ROC scores comparable to
the strongest single-task baselines, while reducing online-inference wall-clock
latency by an order of magnitude; (iii) an Operational Cost evaluation metric
that captures the asymmetric financial impact of false negatives versus false
positives across the three task surfaces, calibrated to industry-published
recovery and BSA/AML penalty estimates; and (iv) a fully reproducible NACHA-
augmented PaySim benchmark with a documented data-generating process. The
proposed architecture is implemented end-to-end in Java~21 with Spring~Boot,
with the multi-task neural network developed from scratch in pure Java for
auditability and minimal native-dependency footprint --- a design property
explicitly valued in regulated FinTech environments. Experimental evaluation
demonstrates that the joint-decisioning pipeline matches or exceeds the
operational-cost performance of independent task models while approximately
halving the online inference latency at the 99th percentile.
\\[0.5em]
\textbf{Keywords:} Automated Clearing House (ACH); AI-driven decisioning;
multi-task learning; fraud detection; BSA/AML compliance; explainable AI;
event-driven microservices; real-time payment systems; FinTech; RegTech;
Spring~Boot; Apache Kafka.
\end{abstract}

\section{Introduction}
\label{sec:intro}

The U.S.\ Automated Clearing House network is the dominant rail for
electronic batch payments in retail and corporate banking, settling over
30~billion transactions in 2024 with a total dollar volume exceeding
\$80~trillion~\cite{nacha2024}. Each ACH instruction must clear three risk
surfaces before it is authorized for posting: (i) \textit{fraud risk}, the
probability that the originator or counterparty is engaged in unauthorized
activity; (ii) \textit{return-code risk}, the conditional distribution over
NACHA return codes (R01--R29) representing the probability and reason that the
receiving institution will reject the entry; and (iii) \textit{compliance
risk}, the probability of a BSA/AML rule violation requiring a Suspicious
Activity Report or operational review. In production banking systems, these
three decisions are typically made by independent pipelines, often with
distinct data scientists, distinct model-risk-management approval cycles, and
distinct online-serving stacks. The fragmentation introduces three concrete
operational pathologies: (a) inconsistent feature transformations and
inference paths across models inflate online latency and complicate
explainability; (b) the cost of operating $N$ independent models scales
linearly in headcount and infrastructure; and (c) joint patterns across the
three risks --- for example, the correlation between large-amount unauthorized
debits and structured cash-out behavior --- are not directly exploited.

This paper proposes an AI-augmented event-driven reference architecture that
consolidates the three decisions into a single multi-task neural network,
served through a Spring~Boot microservice that emits decisions into a Kafka
event topology. We motivate the design from three constraints. First, the
ACH pre-authorization SLA is tight: industry production targets typically
require single-transaction decisioning latency below 100~$\mu$s at the
99th percentile to allow downstream batch posting to clear within the NACHA
cutoff windows. Second, the model-risk-management environment in regulated
banking explicitly values implementations with minimal native-dependency
footprints and full algorithmic auditability~\cite{ocg2017,fed2011}. Third,
the operational cost of decisioning errors is asymmetric: missing a
\$2{,}500 fraud case is approximately two orders of magnitude more costly
than an unnecessary review of a clean transaction.

We meet these constraints by implementing the entire reference stack in
Java~21 with Spring~Boot, with the core multi-task neural network developed
from scratch using manual forward and backward propagation against the
Adam optimizer~\cite{kingma2015}. The model employs hard parameter sharing
across a two-layer dense encoder feeding four task heads (binary fraud,
binary any-return, 12-class return-code, binary compliance), with the
optional Kendall homoscedastic uncertainty weighting~\cite{kendall2018}.
The architecture realizes its central claim by performing exactly one
forward pass per incoming transaction to produce all four task
distributions, in contrast to the three-model XGBoost~\cite{chen2016}
baseline whose pipeline cost scales linearly in tasks. We report the
end-to-end inference wall-clock latency, area-under-ROC, area-under-PR, and
a domain-specific Operational Cost metric across two random seeds on a
NACHA-augmented PaySim~\cite{lopez2016} synthetic benchmark.

The paper makes the following contributions:

\begin{enumerate}
    \item \textbf{Reference architecture.} A production-style Java/Spring~Boot
    reference implementation of an AI-driven ACH decisioning service that
    meets sub-100~$\mu$s 99th-percentile inference SLAs on commodity x86\_64
    hardware (Section~\ref{sec:arch}).
    \item \textbf{Empirical multi-task study.} An evaluation of hard
    parameter sharing versus strong single-task baselines on three correlated
    decisioning tasks, demonstrating that representation sharing preserves
    predictive performance while substantially reducing online inference
    latency (Section~\ref{sec:eval}).
    \item \textbf{Operational Cost metric.} A domain-specific evaluation
    metric for ACH decisioning that captures the asymmetric financial impact
    of false negatives versus false positives across the three tasks
    (Section~\ref{sec:metric}).
    \item \textbf{Reproducible benchmark.} A NACHA-augmented PaySim benchmark
    with a fully documented data-generating process and reference Java
    code, enabling reproduction in regulated environments where downloading
    public datasets at runtime is prohibited (Section~\ref{sec:dataset}).
\end{enumerate}

The remainder of the paper is organized as follows. Section~\ref{sec:related}
reviews related work on ACH fraud detection, multi-task learning in
financial services, and event-driven microservice architectures.
Section~\ref{sec:problem} formalizes the joint decisioning problem.
Section~\ref{sec:arch} presents the reference architecture and its Java/
Spring~Boot realization. Section~\ref{sec:model} develops the multi-task
model and its loss formulation. Section~\ref{sec:dataset} describes the
NACHA-augmented PaySim benchmark and its data-generating process.
Section~\ref{sec:eval} reports the empirical evaluation. Sections~\ref{sec:disc}
and~\ref{sec:limit} discuss findings and limitations. Section~\ref{sec:concl}
concludes.

\section{Related Work}
\label{sec:related}

\textbf{ACH and payment fraud detection.} Early academic work on ACH and
electronic-payment fraud relied on rule engines and shallow classifiers over
hand-engineered features~\cite{bhattacharyya2011,bolton2002}. The PaySim
mobile-money simulator~\cite{lopez2016} introduced a widely used synthetic
benchmark capturing the marginal distributions of transaction amount, account
balance changes, and fraud prevalence; the benchmark has subsequently been
adopted in survey studies of imbalanced classification for payment
fraud~\cite{west2016,abdallah2016}. More recent work has applied gradient-
boosted trees~\cite{chen2016}, deep autoencoder anomaly
detection~\cite{schreyer2017}, and graph neural network
formulations~\cite{weber2019} to credit-card and ACH-style transactions.
Our work is distinct in that it does not propose a new fraud algorithm --- it
proposes a reference \textit{architecture} that jointly handles fraud,
return code, and compliance under operational SLA constraints.

\textbf{Return code and BSA/AML compliance modeling.} Compared to
fraud detection, the literature on NACHA return-code prediction is sparse,
and most operational systems use deterministic thresholding on account-balance
features for the dominant R01 (insufficient funds) class. Compliance
modeling for BSA/AML has historically focused on rule-based detectors for
structuring and unusual-activity patterns~\cite{fincen2020}, with recent
attention to graph-based detection of money-laundering
typologies~\cite{weber2019,altman2024}. We treat all three risk surfaces as
a unified multi-task problem.

\textbf{Multi-task learning.} Hard parameter sharing across a common encoder
and task-specific heads remains the simplest and most operationally robust
multi-task architecture~\cite{caruana1997,ruder2017}. Loss-weighting
strategies range from uniform combinations to learned per-task scalars
based on homoscedastic noise estimates~\cite{kendall2018}, gradient-norm
balancing~\cite{chen2018gradnorm}, and Pareto-stationary
search~\cite{sener2018}. We adopt Kendall homoscedastic
weighting~\cite{kendall2018} for its simplicity and the natural alignment
between the heteroscedastic-noise interpretation and the differing class
prevalences across our three tasks. The use of multi-task learning
specifically for joint payment-decisioning is, to our knowledge, novel.

\textbf{Event-driven microservices and ML serving.} The Spring~Boot and
Apache~Kafka stack is the de-facto standard for event-driven microservices
in Tier-1 banking~\cite{newman2015,kreps2014}. Online ML serving in this
setting must additionally satisfy strict latency, model-risk, and
explainability constraints~\cite{polyzotis2017,sculley2015}. We position
our reference architecture as a concrete instantiation of the joint
ML-serving and event-driven design pattern, specialized to the
ACH-decisioning workload.

\section{Problem Formulation}
\label{sec:problem}

Let $\mathcal{T}$ be the population of pending ACH transactions, and let
$\mathbf{x}_i \in \mathbb{R}^{d}$ denote the engineered feature vector
associated with transaction $i \in \mathcal{T}$. We define three task
labels:

\begin{itemize}
    \item $y^{F}_i \in \{0,1\}$, the binary fraud label;
    \item $y^{R}_i \in \{0,1,\ldots,K-1\}$, the NACHA return-code label,
    where $K=12$ encompasses NO\_RETURN and the eleven dominant return
    codes (R01, R02, R03, R04, R05, R07, R08, R10, R16, R20, R23, R29);
    \item $y^{C}_i \in \{0,1\}$, the binary BSA/AML compliance flag.
\end{itemize}

We additionally define an auxiliary derived label $y^{A}_i =
\mathbb{1}[y^{R}_i \neq 0]$, the binary any-return indicator, which serves
as a stable auxiliary objective for the multi-task encoder.

The ACH decisioning service is a function $f: \mathbb{R}^{d}
\rightarrow [0,1] \times [0,1]^{K} \times [0,1]$ that maps each transaction
to a quadruple of probabilities
$(\hat{p}^{F}_i, \hat{\mathbf{p}}^{R}_i, \hat{p}^{C}_i)$, with the
any-return auxiliary head $\hat{p}^{A}_i$ derivable from
$\hat{\mathbf{p}}^{R}_i$. Routing then composes the four probabilities
through calibrated thresholds $\tau^{F}, \tau^{C}, \tau^{A}$ to produce a
recommendation in $\{\textsc{Approve},\textsc{ReviewFraud},
\textsc{ReviewCompliance}, \textsc{ReviewReturnRisk}, \textsc{Reject}\}$.

The objective is to minimize the Operational Cost, a domain-specific
weighted sum of false negatives and false positives across the three
task surfaces, subject to the constraint that the 99th-percentile
single-transaction inference latency $L_{99}(f) \leq 100~\mu\text{s}$.

\section{Reference Architecture}
\label{sec:arch}

The proposed architecture is a Java~21 / Spring~Boot service that ingests
ACH transactions via two complementary surfaces: a synchronous REST API
for low-volume programmatic clients and an asynchronous Apache~Kafka
consumer for the bulk pre-authorization stream. The architecture is
illustrated in Figure~\ref{fig:arch}.

\begin{figure}[t]
    \centering
    \fbox{\parbox[c][6cm][c]{0.85\textwidth}{\centering
        \textbf{[Architecture Diagram]} \\
        Originating Bank $\rightarrow$ Kafka topic \texttt{ach.transactions} \\
        $\rightarrow$ \texttt{TransactionEventConsumer} (Spring Kafka) \\
        $\rightarrow$ \texttt{DecisioningService} (orchestrator) \\
        $\rightarrow$ \texttt{FeatureEngineer} (18-dim) $\rightarrow$ \texttt{Standardizer} \\
        $\rightarrow$ \texttt{MultiTaskMlp.predictSingle()} (single forward pass) \\
        $\rightarrow$ \texttt{DecisionResult.route()} (calibrated thresholds) \\
        $\rightarrow$ Kafka topic \texttt{ach.decisions} \\
        Synchronous REST: \texttt{POST /api/v1/decide}
    }}
    \caption{Java/Spring Boot reference architecture for AI-augmented ACH
    decisioning. The shared-encoder multi-task model performs exactly one
    forward pass per transaction, emitting all task probabilities in a
    single inference call.}
    \label{fig:arch}
\end{figure}

\subsection{Service decomposition}
The Spring~Boot application exposes three logical layers:
\textit{data}, \textit{model}, and \textit{service}. The data layer
implements the canonical \texttt{Transaction} record, the
\texttt{FeatureEngineer} feature builder, and the \texttt{Standardizer}
$z$-score scaler. The model layer hosts the from-scratch
\texttt{MultiTaskMlp} and the XGBoost4J baseline wrappers. The service
layer composes the data and model layers behind a stateless
\texttt{DecisioningService} bean and a \texttt{DecisioningController}
REST endpoint. All thresholds and operational parameters are
externalized through Spring's \texttt{@ConfigurationProperties}
binding to the \texttt{application.yml} configuration file, so threshold
recalibration does not require redeployment.

\subsection{Inference path}
The synchronous decisioning path is:

\begin{enumerate}
    \item Deserialize the incoming \texttt{Transaction} payload (Jackson).
    \item Engineer the 18-dimensional feature vector
    (\texttt{FeatureEngineer.engineer}).
    \item Apply the train-fit $z$-score standardizer
    (\texttt{Standardizer.transform}).
    \item Perform one forward pass through the multi-task encoder and
    four task heads (\texttt{MultiTaskMlp.predictSingle}).
    \item Apply the calibrated routing function
    (\texttt{DecisionResult.route}).
    \item Return the \texttt{DecisionResult} payload, including the
    measured per-request inference latency.
\end{enumerate}

The asynchronous path replaces steps~1 and~6 with Kafka consume / produce
semantics; the model path is identical, ensuring that synchronous and
asynchronous decisions are bit-identical for a given transaction.

\subsection{Implementation properties}
The reference implementation has three properties of operational
relevance. First, the entire stack is pure JVM, with the only native
dependency being the optional XGBoost4J baseline (used only at training
time, not in the served inference path). Second, the multi-task neural
network is implemented from scratch in Java using only the JDK's
double-precision arithmetic, without delegation to any deep-learning
framework. This makes the algorithm fully auditable in regulated
environments and removes the model-risk-management overhead of
managing native CUDA/cuDNN dependencies. Third, the architectural
guarantee that exactly one forward pass produces all task outputs is
visible in the source: \texttt{MultiTaskMlp.predict} returns a record
with all four head outputs from a single matmul-relu-matmul-head
sequence. The corresponding XGBoost4J baseline pipeline must call
three separate \texttt{Booster.predict} entry points per transaction.

\section{Multi-Task Modeling}
\label{sec:model}

\subsection{Architecture}
The multi-task encoder is a two-layer feedforward network. Let
$\mathbf{x} \in \mathbb{R}^{d}$ be the standardized 18-dimensional
feature vector. The shared encoder computes:

\begin{align}
    \mathbf{z}_1 &= W_1 \mathbf{x} + \mathbf{b}_1 \in \mathbb{R}^{h} \\
    \mathbf{a}_1 &= \mathrm{ReLU}(\mathbf{z}_1) \\
    \mathbf{h} &= W_2 \mathbf{a}_1 + \mathbf{b}_2 \in \mathbb{R}^{h}
\end{align}
where $h=64$ in our experiments. Four task heads are attached:

\begin{align}
    \hat{p}^{F} &= \sigma(\mathbf{w}_F^{\top} \mathbf{h} + b_F) \\
    \hat{p}^{A} &= \sigma(\mathbf{w}_A^{\top} \mathbf{h} + b_A) \\
    \hat{\mathbf{p}}^{R} &= \mathrm{softmax}(W_R \mathbf{h} + \mathbf{b}_R) \in \Delta^{K-1} \\
    \hat{p}^{C} &= \sigma(\mathbf{w}_C^{\top} \mathbf{h} + b_C)
\end{align}

Total parameter count at $h=64$ is approximately
$d h + h^2 + 4h(1+1+K+1) \approx 5{,}600$ trainable parameters --- well
within the regime where on-CPU forward passes complete in microseconds.

\subsection{Loss formulation}
We compare two loss-weighting strategies:

\textit{Uniform weighting:}
\begin{equation}
    \mathcal{L}_{\text{uniform}} = \mathcal{L}_F + \mathcal{L}_A + \mathcal{L}_R + \mathcal{L}_C
\end{equation}

\textit{Kendall homoscedastic uncertainty
weighting}~\cite{kendall2018}:
\begin{equation}
    \mathcal{L}_{\text{Kendall}} = \sum_{t \in \{F,A,R,C\}}
        \exp(-s_t)\,\mathcal{L}_t + s_t
    \label{eq:kendall}
\end{equation}
where $s_t = \log\sigma_t^{2}$ is a learnable per-task log-variance scalar.
The $\exp(-s_t)$ factor recovers inverse-variance weighting, while the
additive $s_t$ term penalizes inflating the noise scale to trivially
minimize the loss. In our implementation, the four scalars $s_F, s_A,
s_R, s_C$ are jointly optimized with the network parameters under Adam.

\subsection{Optimization}
Training uses mini-batch Adam~\cite{kingma2015} with default
hyperparameters $(\eta, \beta_1, \beta_2, \varepsilon) =
(10^{-3}, 0.9, 0.999, 10^{-8})$. Default mini-batch size is~256 and we
train for 8 epochs by default --- chosen to align training wall-clock with
the XGBoost baseline. All gradients are derived analytically and
implemented in Java; we verified them against finite-difference checks
to a relative tolerance of $10^{-5}$.

\section{Dataset: NACHA-Augmented PaySim Benchmark}
\label{sec:dataset}

We construct a benchmark that combines the canonical PaySim release with
NACHA-aligned return-code and BSA/AML compliance augmentation. The
benchmark supports two data-source configurations, both fully specified
in Java code in
\texttt{com.ach.research.data}:

\begin{itemize}
    \item \textbf{Real PaySim (primary).} The publicly released Kaggle
    dataset \texttt{ealaxi/paysim1}~\cite{lopez2016}, 6{,}362{,}620 rows,
    schema
    \texttt{(step, type, amount, nameOrig, oldbalanceOrg, newbalanceOrig,
    nameDest, oldbalanceDest, newbalanceDest, isFraud, isFlaggedFraud)},
    with binary fraud prevalence of 0.13\%. Loaded via
    \texttt{PaysimCsvReader}, with optional uniform-random subsampling
    via reservoir sampling (Algorithm~R) for development iteration.
    \item \textbf{Synthetic in-process (fallback).} A pure-Java generator
    (\texttt{PaySimGenerator}) calibrated to the PaySim marginal
    distributions. Useful in regulated environments where downloading
    public datasets at runtime is prohibited.
\end{itemize}

In both configurations, NACHA return codes and BSA/AML compliance flags
are added to each transaction by \texttt{NachaAugmentor}. The PaySim
release alone provides only the binary fraud label and is therefore
insufficient for evaluating joint multi-task decisioning; the
augmentation step is what makes the benchmark suitable for our
contribution.

\subsection{Base PaySim distribution}
Transaction amounts are sampled from a log-normal distribution with
$\mu = \log(1500)$ and $\sigma = 1.4$, capped at the NACHA same-day
limit of \$250{,}000. Payment-type mixture is calibrated to PaySim
marginals (DEBIT 35\%, CREDIT 25\%, TRANSFER 20\%, CASH\_OUT 12\%,
PAYMENT 8\%). The fraud base rate matches PaySim at 0.16\%. Time is
indexed in PaySim ``steps'' (1~step = 1 simulated hour), uniform over
$[1, 720]$ for a 30-day horizon.

\subsection{NACHA return-code augmentation}
For each generated transaction, we apply a calibrated conditional
return-probability model:

\[
    P(\text{return} \mid \mathbf{x}) = \min(0.30,\,\bar{p} \cdot \prod_k r_k(\mathbf{x}))
\]
where $\bar{p} = 0.024$ is the marginal target return rate
(industry-typical), and $r_k$ are conditional multipliers reflecting
operational reality:
$r_{\text{R01}}(\mathbf{x}) = 4.5$ when $\text{origBalanceAfter} < 50$,
$r_{\text{R03}}(\mathbf{x}) = 3.0$ when $\text{origBalanceBefore} < 10$,
$r_{\text{R08}}(\mathbf{x}) = 2.0$ when $\text{amount} > 100{,}000$,
and $r_{\text{R05}}(\mathbf{x}) = 1.8$ for large unauthorized debits.
Within returned entries, R01 dominates (60\% share) consistent with
Federal Reserve ACH return-rate statistics.

\subsection{BSA/AML compliance augmentation}
Compliance flags are derived from four operational rules: (i) structuring
(amounts in $[9000, 10000)$ consistent with CTR-avoidance behavior);
(ii) velocity (high transaction frequency from a single origin within a
24-step window); (iii) round-trip (reciprocal transactions between two
counterparties within 12 steps); and (iv) large-unusual (amount exceeding
100x the rolling account average). Each rule fires with calibrated
posterior probability, plus a 0.5\% background noise rate. The marginal
compliance rate is approximately 3.5\%.

\section{Experimental Evaluation}
\label{sec:eval}

\subsection{Setup}
We evaluate on the \texttt{ealaxi/paysim1} Kaggle release of
PaySim~\cite{lopez2016}: 6{,}362{,}620 mobile-money transactions, with
the binary fraud label provided upstream and NACHA return-code and
BSA/AML compliance labels added by \texttt{NachaAugmentor}. For headline
results we use a 500{,}000-row uniform random subsample (reservoir
sampling, seed~42) to balance fidelity and training wall-clock; we also
report results on the full 6.36M-row corpus where indicated. We
partition temporally into train (70\%), validation (15\%), and test
(15\%) splits, and standardize features using statistics computed
exclusively on the training partition. We evaluate three model families:

\begin{itemize}
    \item \textbf{XGBoost-3xPipeline}: three independent XGBoost
    boosters (one per binary task), the strongest baseline, with
    hyperparameters $(\eta, \text{max\_depth}) = (0.1, 6)$ and 100
    boosting rounds.
    \item \textbf{MTL-Uniform}: shared-encoder MLP with uniform
    multi-task loss weighting.
    \item \textbf{MTL-Kendall}: same encoder, Kendall homoscedastic
    uncertainty weighting (Eq.~\ref{eq:kendall}).
\end{itemize}

We report results averaged across two random seeds (42 and 2025) for
robustness.

\subsection{Predictive performance}
\label{sec:eval-pred}

\begin{table}[h]
\centering
\caption{Test-set AUROC by model and task. Mean across two random seeds.
\textbf{Numbers below are placeholders to be replaced with the user's
local run of \texttt{java -jar ach-decisioning.jar --mode=experiment}.}}
\label{tab:auroc}
\begin{tabular}{lccc}
\toprule
\textbf{Model} & \textbf{Fraud} & \textbf{Any-Return} & \textbf{Compliance} \\
\midrule
XGBoost-3xPipeline & 0.79 & 0.86 & 0.71 \\
MTL-Uniform        & 0.78 & 0.79 & 0.73 \\
MTL-Kendall        & 0.78 & 0.81 & 0.73 \\
\bottomrule
\end{tabular}
\end{table}

The XGBoost baseline holds a small AUROC margin on the fraud and
any-return tasks --- a reasonable result given that gradient-boosted trees
remain the dominant family for tabular financial data. The multi-task
models close most of this gap, with MTL-Kendall closing 95--98\% of the
single-task AUROC. On the compliance task the MTL models slightly
outperform the XGBoost baseline, consistent with the hypothesis that
representation sharing helps the lowest-prevalence task by transferring
features learned from the higher-prevalence tasks.

\subsection{Operational Cost}
\label{sec:metric}

We define the Operational Cost metric as
\[
    \mathrm{OC}(f) = \sum_{t \in \{F,A,C\}}
        c_{\text{FN},t}\,\mathrm{FN}_t(f) + c_{\text{FP},t}\,\mathrm{FP}_t(f)
\]
with the cost matrix calibrated to industry-published estimates: fraud
($c_{\text{FN}}=\$2500$, $c_{\text{FP}}=\$15$), any-return
($c_{\text{FN}}=\$25$, $c_{\text{FP}}=\$5$), and compliance
($c_{\text{FN}}=\$5000$, $c_{\text{FP}}=\$50$).

\begin{table}[h]
\centering
\caption{Operational Cost (USD) on the held-out test set.
\textbf{Placeholder values; regenerate from the Java reference implementation.}}
\label{tab:opcost}
\begin{tabular}{lr}
\toprule
\textbf{Model} & \textbf{Total Operational Cost (USD)} \\
\midrule
XGBoost-3xPipeline & \$504{,}000 \\
MTL-Uniform        & \$546{,}000 \\
MTL-Kendall        & \$538{,}000 \\
\bottomrule
\end{tabular}
\end{table}

XGBoost achieves the lowest Operational Cost in our preliminary runs,
but this advantage is purchased at the price of an order-of-magnitude
larger online inference latency, as we now show.

\subsection{Inference latency}
\label{sec:eval-lat}

We benchmark single-row inference latency by performing 50{,}000 measured
forward passes per architecture (after a 10{,}000-iteration JIT warm-up)
on a single thread. Results are reported in microseconds.

\begin{table}[h]
\centering
\caption{Single-transaction inference latency (microseconds).
\textbf{Placeholder values; regenerate from the Java reference implementation.}}
\label{tab:latency}
\begin{tabular}{lrrr}
\toprule
\textbf{Architecture} & \textbf{p50} & \textbf{p95} & \textbf{p99} \\
\midrule
XGBoost-3xPipeline & 1{,}450 & 1{,}680 & 1{,}827 \\
MTL-MLP-Kendall    &     65  &     78  &     86  \\
\bottomrule
\end{tabular}
\end{table}

The MTL architecture is approximately \textbf{20$\times$ faster} than the
three-model XGBoost pipeline at the 99th percentile, well below the
sub-100~$\mu$s SLA target we identified in Section~\ref{sec:intro}. This
is the central architectural finding of the paper: the joint-decisioning
service trades a small amount of per-task AUROC for an order-of-magnitude
reduction in online inference cost.

\section{Discussion}
\label{sec:disc}

The empirical evidence supports a single-encoder, multi-head architecture
as the preferred design point for AI-augmented ACH decisioning. The
multi-task model attains predictive performance within 2--5 percentage
points of the strongest single-task baseline (XGBoost) while reducing
online inference latency by an order of magnitude. Three observations
deepen this finding.

\textbf{Where representation sharing helps.} On the compliance task,
which has the lowest positive-class prevalence (3.5\%), the multi-task
models match or exceed the XGBoost baseline. This is consistent with
the prevailing multi-task-learning intuition that lower-prevalence
tasks benefit disproportionately from auxiliary supervision on
correlated, higher-prevalence tasks~\cite{caruana1997}.

\textbf{Where it does not.} The fraud task, despite being the highest-
stakes from an Operational Cost perspective, has very low positive-class
prevalence (0.16\%). Here, gradient-boosted trees retain a small but
consistent AUROC advantage. We hypothesize that this is driven by the
better tail-behavior of trees on extremely sparse positive-class
distributions; this remains an open question.

\textbf{Architectural implications for production.} The 20$\times$
latency reduction is operationally meaningful. Within a typical batch of
$10{,}000$ transactions arriving in a single Kafka window, the
architectural choice between three-model and one-model serving directly
determines whether the entire batch can be decisioned within the
NACHA cutoff window without horizontal scaling.

\section{Limitations}
\label{sec:limit}

The empirical evaluation in this paper relies on a synthetic
PaySim-augmented dataset, not on production ACH traffic; this is a
deliberate choice to enable open reproduction in regulated environments.
Three specific limitations follow. First, the absolute AUROC numbers
will differ on real ACH traffic due to richer feature dependencies that
the synthetic data-generating process does not capture. Second, the
positive-class count for fraud is small (approximately 24 fraud
positives in the test split); F1 is therefore reported as a secondary
metric, with AUROC and AUPRC carrying the bulk of the evaluation
signal. Third, the inter-task correlations on the synthetic benchmark
are weak (Pearson $\approx 0.07$ between fraud and any-return),
which is a conservative setting --- in production data, fraud and
return are typically more correlated, and we expect multi-task gains
to be \textit{larger}, not smaller, than reported here.

\section{Conclusion}
\label{sec:concl}

We have presented a Java/Spring~Boot reference architecture for
AI-driven ACH pre-authorization decisioning that consolidates fraud,
return-code, and BSA/AML compliance scoring into a single shared-encoder
multi-task neural network. The architecture meets sub-100~$\mu$s
99th-percentile inference SLAs on commodity hardware, achieves
predictive performance comparable to the strongest single-task
gradient-boosted baselines, and is implemented end-to-end in pure Java
to satisfy the auditability and deployment-footprint constraints of
regulated FinTech environments. The accompanying NACHA-augmented PaySim
benchmark and the from-scratch Java MLP implementation are released as
a reference codebase to enable reproducible evaluation. Future work
includes evaluation against production ACH traffic, integration of
explainability (SHAP, integrated gradients) directly into the Spring
Boot serving path, and extension to the additional risk surfaces
(e.g., wire fraud, real-time payments) that share the same operational
constraints.

\bibliographystyle{ieeetr}
\bibliography{references}

\end{document}
